\documentclass[12pt]{article}
\usepackage[utf8]{inputenc}
\usepackage[T1]{fontenc}
\usepackage{hyperref}
\title{Tatana Language Resources for K-12 Curriculum Development}
\author{Compiled by Claude Sonnet 5 (claude-sonnet-5) for Andr\'as Kornai}
\date{}
\begin{document}
\maketitle
\section{Introduction}
Tatana' (ISO 639-3 txx) is an Austronesian language of Sabah, Malaysia, spoken in Kuala Penyu, Kota Kinabalu and the Klang Valley. Its classification is uncertain: it may be a Bisaya variety, and Lobel (2013) places it in Murutic. Searches found no Tatana-specific school materials.
\section{Primary Educational Resources}
None found in search results.
\section{Linguistic and Reference Materials}
\begin{itemize}
\item Tatana language, Wikipedia overview. It lists a trilingual Tatana'--Malay--English phrase book (Pekkanen, Dunn Chan and Dillon 1998, Sabah Museum), which was not otherwise located: \url{https://en.wikipedia.org/wiki/Tatana_language}
\item Tatana numerals list (Max Planck, from John Dillon, SIL): \url{https://lingweb.eva.mpg.de/channumerals/Tatana.htm}
\item Glottolog reference record: \url{https://glottolog.org/resource/reference/id/132711}
\end{itemize}
\section{Religious and Community Resources}
None found in search results.
\section{Digital and Audio Resources}
None found in search results.
\section{Curriculum Development Context}
Very little is documented. The phrase book is the closest thing to a teaching resource. Related Kadazan-Dusun and Bisaya materials might serve as comparative background, but none were verified here.
\end{document}